\subsection{Паттерны по типам запросов и архитектурные наблюдения}

Объединение результатов двух серий позволяет выделить устойчивые закономерности
(таблица~\ref{tab:patterns}). В потоковом сценарии сравнение ведётся по медиане как
менее зашумлённой метрике, отражающей типичную задержку, к которой и чувствителен
потребитель витрины данных.

\begin{table}
    \caption{Сравнение форматов по типам запросов и операций (NRT-сценарий, медиана)}
    \label{tab:patterns}
    \begin{tabularx}{\textwidth}{|>{\raggedright\arraybackslash}p{0.24\textwidth}|>{\raggedright\arraybackslash}p{0.20\textwidth}|>{\raggedright\arraybackslash}p{0.20\textwidth}|X|}\hline
    \textbf{Тип операции} & \textbf{Vortex против Parquet} & \textbf{Vortex против ORC} & \textbf{Объяснение} \\ \hline
    Запись (INIT) & сопоставимы & сопоставимы & оба буферизуют строки и сбрасывают по чекпойнту, различия в пределах разброса \\ \hline
    Точечный доступ по ключу (\texttt{POINT}) & сопоставимы (разница доли процента) & сопоставимы & фильтры Блума и статистики работают у всех; на маленьких выборках различия в пределах разброса \\ \hline
    Полное сканирование с проекцией (\texttt{FULLSCAN\_PROJ}) & сопоставимы (разница доли процента) & Vortex чуть быстрее в FINAL ($\approx 1\,\%$) & column pruning + плотные кодеки работают на упорядоченных данных; на свежих L0-файлах эффект слабее \\ \hline
    Диапазонный фильтр по строке-префиксу (\texttt{RANGE\_SESSION}, \texttt{RANGE\_URL}) & Vortex заметно быстрее (\texttt{RANGE\_SESSION} в DURING $\approx 3{\times}$ быстрее) & Vortex стабильно быстрее (\texttt{RANGE\_SESSION} в DURING $\approx 25\,\%$; \texttt{RANGE\_URL} $\approx 4\,\%$) & FSST-кодирование строк с общими префиксами + нативное проталкивание предиката + пропуск файлов по полным статистикам \\ \hline
    Подсчёт с селективным предикатом (\texttt{AGG\_FILTERED}) & Vortex быстрее ($\approx 21\,\%$ в FINAL) & Vortex стабильно быстрее ($\approx 21\,\%$ в FINAL, $\approx 4\,\%$ в DURING) & предикат отсекает данные до пересечения границы JNI; выигрыш от проталкивания преобладает над накладными расходами \\ \hline
    Многошаговая аналитика (\texttt{FUNNEL}, тяжёлые соединения с фильтром) & Vortex стабильно быстрее (\texttt{FUNNEL} $\approx 33\,\%$) & Vortex стабильно быстрее (\texttt{FUNNEL} $\approx 33\,\%$, \texttt{JOIN\_REVENUE\_PROFILE} $\approx 20\,\%$) & многократные проходы и предикаты, выигрыш от column pruning и проталкивания умножается \\ \hline
    Простая группировка/многомерная агрегация без предиката (\texttt{AGG\_GROUP\_COUNTRY}, \texttt{AGG\_MULTIDIM}) & Parquet иногда быстрее (единицы процентов) & ORC/Parquet иногда быстрее (единицы процентов) & чистая агрегация в JVM, для Vortex --- через границу JNI; нет проталкивания агрегатов в интерфейсе Paimon \\ \hline
    Хвосты задержки под конкурентной записью (среднее против медианы) & Vortex существенно у́же (среднее $\approx$ медиана, у Parquet среднее на $\approx 21\,\%$ выше медианы) & оба формата дают узкие хвосты (среднее $\approx$ медиана) разными способами & пул соединений Hadoop S3A эпизодически насыщается на сбросе чекпойнта; Vortex изолирован независимым нативным пулом, ORC --- агрессивной буферизацией C++-читателя \\ \hline
    \end{tabularx}
\end{table}

\textbf{Где Vortex реально лучше ORC.} Если выбрать устойчивые наблюдения и
оставить запросы, на которых разница превышает разброс, преимущество Vortex над ORC
сосредоточено в трёх классах запросов: запросы с селективным предикатом
(проталкивание в нативный декодер до границы JNI), диапазонные запросы по строкам с
регулярной структурой (FSST + пропуск по полным статистикам), многошаговая аналитика
с несколькими предикатами и проходами. На простой группировке по столбцу низкой
кардинальности и на многомерной агрегации без предиката ORC оказывается не хуже
или лучше, на точечных запросах --- сопоставим.

\textbf{Эффект упорядоченности данных.} По медиане ORC и Parquet практически не
меняют времени между фазами FINAL и DURING (отношение DURING/FINAL близко к единице),
а Vortex в DURING медленнее самого себя в FINAL на $\approx 22\,\%$. Это
прямое следствие зависимости специализированных кодеков Vortex (прежде всего FSST)
от \textit{порядка} строк внутри блока: в свежих файлах уровня~0, формируемых
по чекпойнту, данные ещё не отсортированы по ключу слиянием, и плотность
кодирования снижается; после остановки записи и работы фонового слияния файлы
сортированы по ключу, и кодеки работают в полную плотность. ORC и Parquet таким
эффектом не обладают --- их кодеки менее чувствительны к порядку, --- поэтому их
типичная задержка между фазами и не меняется.

\textbf{О выборе метрики усреднения и хвостах задержки.} В DURING-фазе отдельные
сэмплы чтения могут попадать в момент массового сброса данных по чекпойнту, когда
writer одновременно пишет в объектное хранилище и данные таблицы, и состояние
чекпойнта Flink (последнее размещается в той же директории объектного хранилища).
Среднее по сэмплам, особенно при небольшом числе повторов, чувствительно к таким
попаданиям. В работе в качестве основной метрики использована медиана. Соотнесение
среднего и медианы по DURING-фазе показывает: у ORC и Vortex среднее и медиана
практически совпадают (узкие хвосты задержки), у Parquet среднее на $\approx 21\,\%$
выше медианы --- среднее инфлировано несколькими крупными выбросами в десятки
секунд. Это согласуется с архитектурным предположением о пуле соединений: читатель
Parquet ходит в объектное хранилище через общий пул слоя совместимости Hadoop S3A
и в момент сброса данных по чекпойнту иногда простаивает в ожидании соединения, что
и даёт катастрофически медленную выборку. Читатель Vortex использует независимый
нативный пул и таких выбросов не даёт. Читатель ORC, формально использующий тот же
общий пул, в этом эксперименте также не даёт выбросов --- за счёт агрессивной
буферизации ввода-вывода в нативном C++-читателе, поглощающей эпизодические
задержки S3-операций.

\textbf{Почему Vortex доминирует на TPC-DS, но не доминирует на потоковом
LSM-сценарии.} На TPC-DS таблицы только для добавления, файлы крупные, путь чтения
прямой колоночный, накладные расходы границы JNI амортизируются на больших пакетах,
а запросы --- с высокой селективностью. На потоковом сценарии таблица --- с первичным
ключом и слиянием при чтении: итератор слияния открывает много файлов (особенно
свежих, не слитых файлов уровня~0), накладные расходы границы JNI накапливаются на
пер-файловой основе, данные в свежих файлах не упорядочены по ключу --- кодеки
Vortex (прежде всего FSST) теряют часть эффективности, --- а селективность запросов
в среднем ниже. Поэтому преимущество Vortex здесь скромнее и проявляется в первую
очередь на упорядоченном итоговом снимке и на запросах с селективным предикатом.
Для Vortex \textit{область наибольшей применимости} --- именно batch-аналитика на
таблицах только для добавления, не LSM с первичным ключом.

\textbf{Почему Parquet иногда быстрее на агрегации при полном сканировании.}
Проталкивание агрегатов отсутствует в интерфейсе файловых форматов Apache Paimon в
принципе --- это не специфика Vortex. Агрегаты вычисляются движком Flink над
декодированными данными. Для Parquet весь конвейер --- декодер и агрегатор --- в JVM,
границы нет. Для Vortex данные декодируются нативно, затем пересекают границу JNI
(без копирования, но с фиксированной стоимостью на пакет) и агрегируются в JVM; на
миллионах пакетов это даёт заметную добавку. Устранение ограничения --- проталкивание
агрегатов в нативный SIMD-слой Vortex --- естественное направление развития Apache
Paimon и потенциальный вклад в его экосистему.

\textbf{Почему два менеджера задач на узел лучше одного крупного.} Меньший размер
кучи на процесс ускоряет полную сборку мусора и сокращает паузы; изоляция процессов
означает, что пауза одного менеджера задач не блокирует второй; это особенно важно
для Parquet, чей путь декодирования и агрегации полностью в куче JVM, при высоком
параллелизме. Компактор Vortex использует память вне кучи и меньше зависит от размера
кучи, но также выигрывает от изоляции.

\subsection{Рекомендации по применению форматов хранения}

Рекомендации сформулированы с учётом сценария использования.

\textbf{1. Batch-аналитика, таблицы только для добавления.} Это \textit{основная
область применимости Vortex}. На наборе TPC-DS масштаба \SI{100}{\giga\byte} Vortex
быстрее Parquet на \num{82} запросах из \num{103}, среднее ускорение --- порядка
\SI{25}{\percent}, пиковое --- до \SI{80}{\percent} (запросы с селективным
предикатом, узкой проекцией, многократным сканированием немногих столбцов). Условия
этого сценария --- крупные файлы, отсутствие слияния версий, упорядоченные данные,
высокая селективность --- максимально благоприятны для специализированных кодеков и
нативного проталкивания предиката. \textbf{Рекомендуется выбирать Vortex} как
формат хранения данных для аналитических витрин на основе таблиц только для
добавления (фактовые таблицы, журналы, исторические данные, материализованные
выборки).

\textbf{2. Near-real-time-витрина (таблица с первичным ключом, потоковая запись).}
По типичной задержке (медиана) \textbf{Vortex} --- быстрейший формат в обеих фазах
и на большинстве запросов (девять из тринадцати); рекомендуется, если в нагрузке
преобладают запросы с селективным предикатом (отсечение до агрегации), диапазонные
запросы по строкам с регулярной структурой (идентификаторы, пути, URL) и
многошаговая аналитика (последовательности событий, соединения с диапазонными
фильтрами). \textbf{ORC} --- разумный выбор, если в нагрузке преобладают простая
группировка по столбцу низкой кардинальности, чистая многомерная агрегация без
предиката или если важна узкая задержка отдельных сэмплов (наихудший случай, p99) ---
ORC демонстрирует одинаково узкие хвосты с Vortex, а в отдельных запросах ему
способствует словарное кодирование. ORC также даёт самый компактный объём хранения
(\SI{17}{\giga\byte} против \SI{18}{\giga\byte} у Parquet и \SI{20}{\giga\byte} у
Vortex на потоковом сценарии; на TPC-DS преимущество ORC меньше). \textbf{Parquet}
--- разумный универсальный выбор по умолчанию: не уступает по записи, хуже по
медиане чтения, зрелый и широко поддержанный; основной недостаток --- эпизодические
выбросы времени запросов под конкурентной записью в моменты сброса данных по
чекпойнту (среднее заметно выше медианы).

\textbf{3. Чего не рекомендуется.} \textbf{Lance} в текущем состоянии интеграции с
Apache Paimon не рекомендуется для потоковых таблиц с первичным ключом (требуется
доработка слияния версий; на OLAP-сценарии отстаёт от остальных форматов в 2--3
раза); его область --- нагрузки машинного обучения с произвольным доступом к строкам
и векторным поиском, выходящие за рамки настоящей работы. \textbf{Apache Avro} ---
строчный формат, в принципе непригоден для OLAP-нагрузки; применение его для
свежих файлов уровня~0 в LSM-дереве было проверено и отклонено --- аналитические
запросы Vortex на итоговом снимке замедляются за счёт совмещения построчных дельт
с колоночными файлами в итераторе слияния.

\textbf{4. Конфигурация (общие соображения для NRT-витрин).} Распределять слоты
между записью и чтением так, чтобы читатель не блокировался; увеличить размер
пакета записи Vortex до десятков тысяч строк (снижает накладные расходы JNI при
чтении в десятки раз); включить полные статистики по столбцам
(\texttt{metadata.stats-mode = full}); настроить расширенные политики хранения
снимков и манифестов под частые фиксации; не применять агрессивных тонких
настроек размера файлов и страниц (благоприятствуют Parquet и раздувают
хранилище) и не использовать строчный формат для свежих файлов уровня~0;
увеличивать интервал чекпойнтов следует с учётом того, что это укрупняет всплеск
ввода-вывода при сбросе.
